% =============================================================================
% ANNEX C — PRIORITISED BACKLOG EXTRACT
% =============================================================================

\chapter{Prioritised Backlog Extract}
\label{annex:backlog}

% --- Note de mise en contexte (portfolio) -------------------------------------
\begin{tcolorbox}[colback=blue!5, colframe=blue!40!black, boxrule=0.5pt, arc=3pt,
  left=8pt, right=8pt, top=6pt, bottom=6pt]
  \small\itshape
  \textbf{Note du portfolio.} Ce document est un extrait du backlog de
  l'\'equipe logicielle TWIICE, export\'e depuis le board GitLab du projet.
  Il illustre la hi\'erarchisation des t\^aches apr\`es application de la
  matrice d'Eisenhower (cf. Phase~2, Espace modules du pr\'esent portfolio).

  Conform\'ement aux r\`egles de confidentialit\'e de l'entreprise, les
  informations sensibles (num\'eros d'issues, noms des collaborateurs, dates,
  jalons internes) ont \'et\'e \textbf{caviard\'ees} et sont repr\'esent\'ees
  par des barres noires (\caviarder[1.2cm]). Les intitul\'es des t\^aches et
  leur description technique sont, quant \`a eux, int\'egralement reproduits.
\end{tcolorbox}

\bigskip

% --- Sprint / Project header --------------------------------------------------
\begin{tcolorbox}[colback=gray!8, colframe=gray!50, boxrule=0.5pt, arc=3pt,
  left=8pt, right=8pt, top=6pt, bottom=6pt]
  \begin{tabularx}{\linewidth}{@{}p{4cm}X@{}}
    \textbf{Project}   & TWIICE Exoskeleton --- Software \\[3pt]
    \textbf{Milestone} & \caviarder[3cm] \\[3pt]
    \textbf{Sprint}    & \caviarder[2cm]\quad
                         (\caviarder[2cm]\;--\;\caviarder[2cm]) \\[3pt]
    \textbf{Board}     & \caviarder[5cm] \\[3pt]
    \textbf{Exported by} & Alexandre Gati \\[3pt]
    \textbf{Export date} & \caviarder[2.5cm] \\
  \end{tabularx}
\end{tcolorbox}

\bigskip

% ==============================================================================
\section*{Context}
% ==============================================================================

The following extract lists the software issues identified and re-prioritised
following an Eisenhower matrix exercise conducted with the team (see Annex~B,
\S1). Issues are grouped by priority level --- \highprio, \mediumprio,
\lowprio{} --- and ordered within each group by creation date. Redacted fields
(\caviarder[0.9cm]) contain information covered by the company's
confidentiality policy.

\bigskip

% ==============================================================================
% BACKLOG TABLE
% ==============================================================================

\begingroup
\setlength{\tabcolsep}{3pt}
\begin{table}[h!]
  \centering\footnotesize
  \begin{tabularx}{\textwidth}{|p{1.4cm}|p{1.1cm}|X|p{2.0cm}|p{1.4cm}|p{1.7cm}|}
    \hline
    \rowcolor{gray!15}
    \textbf{Priority} & \textbf{Issue\,\#} & \textbf{Title \& description} &
    \textbf{Assignee} & \textbf{Due} & \textbf{Status} \\
    \hline

    % ------------------------------------------------------------------
    % HIGH PRIORITY
    % ------------------------------------------------------------------

    \highprio
      & \caviarder[0.8cm]
      & \textbf{Set up SOUP traceability --- GitLab Package Registry}
        \newline{\footnotesize
          Register all third-party libraries (SOUP list) in the GitLab Package
          Registry. Ensure each entry includes version, licence, source URL and
          checksum, per IEC~62304 \S8.1.2.
          \newline\textit{Labels:} \texttt{iec-62304}, \texttt{certification},
          \texttt{sprint-\caviarder[0.4cm]}}
      & \caviarder[2.2cm]
      & \caviarder[1.5cm]
      & \sinprogress \\
    \hline

    \highprio
      & \caviarder[0.8cm]
      & \textbf{Document all third-party components (IEC~62304 \S8)}
        \newline{\footnotesize
          Produce the SOUP list artefact required by IEC~62304 clause~8:
          component name, version, publisher, risk classification.
          Cross-reference with the software requirements specification.
          \newline\textit{Labels:} \texttt{iec-62304}, \texttt{documentation},
          \texttt{sprint-\caviarder[0.4cm]}}
      & \caviarder[2.2cm]
      & \caviarder[1.5cm]
      & \sinprogress \\
    \hline

    \highprio
      & \caviarder[0.8cm]
      & \textbf{Define CI/CD pipeline --- automated build \& test artefacts}
        \newline{\footnotesize
          Configure the GitLab CI pipeline to produce reproducible build
          artefacts (binary, SBOM, test report) and publish them to the
          Package Registry with immutable versioning tags.
          \newline\textit{Labels:} \texttt{ci-cd}, \texttt{iec-62304},
          \texttt{sprint-\caviarder[0.4cm]}}
      & \caviarder[2.2cm]
      & \caviarder[1.5cm]
      & \stodo \\
    \hline

    % ------------------------------------------------------------------
    % MEDIUM PRIORITY
    % ------------------------------------------------------------------

    \mediumprio
      & \caviarder[0.8cm]
      & \textbf{Standardise documentation templates (IEC~62304)}
        \newline{\footnotesize
          Create shared GitLab Wiki templates for: SOUP list, change request,
          software release note, design review record. Reduces manual overhead
          and ensures consistent artefact format across sprints.
          \newline\textit{Labels:} \texttt{documentation},
          \texttt{sprint-\caviarder[0.4cm]}}
      & \caviarder[2.2cm]
      & \caviarder[1.5cm]
      & \stodo \\
    \hline

    \mediumprio
      & \caviarder[0.8cm]
      & \textbf{Update SOUP risk assessment --- IEC~62304 \S7.2}
        \newline{\footnotesize
          Review existing risk assessments for embedded SOUP (FreeRTOS,
          libmodbus, \caviarder[1.5cm]). Align severity classifications with
          the updated device risk file (MDR class~IIa context).
          \newline\textit{Labels:} \texttt{risk}, \texttt{iec-62304}}
      & \caviarder[2.2cm]
      & \caviarder[1.5cm]
      & \sbacklog \\
    \hline

    % ------------------------------------------------------------------
    % LOW PRIORITY
    % ------------------------------------------------------------------

    \lowprio
      & \caviarder[0.8cm]
      & \textbf{Evaluate MCU vs Jetson for user interface module}
        \newline{\footnotesize
          Comparative benchmark: processing power, energy budget, SDK
          maturity, lead time. Decision deferred --- not blocking
          certification. To be scheduled post-milestone
          \caviarder[1.2cm].
          \newline\textit{Labels:} \texttt{hardware}, \texttt{research}}
      & \caviarder[2.2cm]
      & \caviarder[1.5cm]
      & \sbacklog \\
    \hline

    \lowprio
      & \caviarder[0.8cm]
      & \textbf{Benchmark supplementary DevOps tools}
        \newline{\footnotesize
          Evaluate JFrog Artifactory, Nexus Repository and Harbor as
          potential complements to the GitLab Package Registry. No
          blocker identified --- team comfort improvement only.
          \newline\textit{Labels:} \texttt{devops}, \texttt{tooling}}
      & \caviarder[2.2cm]
      & \caviarder[1.5cm]
      & \sbacklog \\
    \hline

  \end{tabularx}
  \caption{Prioritised backlog extract --- TWIICE Software, post-Eisenhower
    matrix triage}
  \label{tab:backlog}
\end{table}
\endgroup

% ==============================================================================
\section*{Traceability to the Eisenhower matrix}
% ==============================================================================

The table below maps each priority group to the corresponding quadrant of the
Eisenhower matrix applied during the Phase~2 (Definition) workshop.

\begin{table}[h!]
  \centering\small
  \begin{tabularx}{\textwidth}{|p{1.6cm}|p{3.5cm}|X|}
    \hline
    \textbf{Priority} & \textbf{Eisenhower quadrant} & \textbf{Rationale} \\
    \hline
    \highprio
      & Urgent \& Important
      & Directly blocks the IEC~62304 certification audit. Must be resolved
        within the current or immediately following sprint. \\
    \hline
    \mediumprio
      & Important, not urgent
      & Required for long-term compliance and team efficiency, but does not
        block the next certification milestone. Scheduled in upcoming sprints. \\
    \hline
    \lowprio
      & Not important (current context)
      & Debated actively during the workshop but identified as non-blocking.
        Deferred to the post-certification backlog. \\
    \hline
  \end{tabularx}
  \caption{Mapping between backlog priority levels and Eisenhower quadrants}
  \label{tab:backlog-eisenhower}
\end{table}





